\documentclass{article}
\usepackage[UTF8]{ctex}
\usepackage{amsmath,amssymb,amsfonts}
\usepackage{enumitem}
\usepackage{geometry}
\geometry{a4paper, margin=1in}
% === 错题本排版优化（不参与题目标识） ===
\setlength{\parindent}{0pt}
\setlength{\parskip}{0.35em plus 0.10em minus 0.08em}
\setlength{\emergencystretch}{3em}
\setlist{itemsep=0.25em, topsep=0.40em, parsep=0pt}
\allowdisplaybreaks[2]
% === 错题本排版优化结束 ===
\begin{document}
\section*{中国矿业大学2026年数学分析真题}
\section*{24．中国矿业大学2026年数学分析试卷}
\begin{itemize}

\subsection{题目1}
  \item[1．] 当 \(x \rightarrow-\infty\) 时，\(f(x)\) 不趋于 \(a\) 的定义为（）。
  \item[A．] \(\exists M>0\) ，对 \(\forall \varepsilon>0, \forall x<-M,|f(x)-a| \geq \varepsilon\) ．
  \item[B．] \(\exists M>0\) ，对 \(\forall \varepsilon>0, \exists x_{0}<-M,\left|f\left(x_{0}\right)-a\right| \geq \varepsilon\) ．
  \item[C．] \(\exists \varepsilon_{0}>0\) ，对 \(\forall M>0, \exists x_{0}<-M,|f(x)-a|<\varepsilon_{0}\) ．
  \item[D．] \(\exists \varepsilon_{0}>0\) ，对 \(\forall M>0, \exists x_{0}<-M,|f(x)-a| \geq \varepsilon_{0}\) ．

\subsection{题目2}
  \item[2．] \(f(x)=\left\{\begin{array}{ll}x+2, & x \geq 0, \\ x-2, & x<0 .\end{array}\right.\) 则 \(f(x)\) 在 \(x=0\) 处（ ）
  \item[A．] 连续．B．左连续但不右连续．C．右连续但不左连续．D．既不左连续也不右连续．
\end{itemize}

\subsection{题目3}
3．变限积分 \(f(x)=\int_{0}^{x} t e^{t} \mathrm{~d} t\) ，则 \(f(x)=()\)

\begin{itemize}
  \item[A．] \(x e^{x}+e^{x}\) ．B．\(x e^{x}+e^{x}\) ．

\subsection{题目4}
\end{itemize}

\subsection{题目4}
原题边界“4．”：若级数 \(\sum_{n=1}^{\infty} a_{n}\) 条件收敛，则 \(\sum_{n=1}^{\infty}\left(a_{n}+\frac{1}{n^{2}}\right)()\)
  \item[A．] 条件收敛．B．绝对收敛．C．发散．D．可能条件收敛也可能绝对收敛．

\subsection{题目5}
5．积分 \(\oint_{L} \frac{x \mathrm{~d} y-y \mathrm{~d} x}{x^{2}+y^{2}}=()\) ，其中 \(L\) 为 \(x^{2}+2 y^{2}=1\) ，取逆时针方向。

\begin{itemize}
  \item[A．] 0．B．\(\pi\) ．C． \(2 \pi\) ．D．\(-\pi\) ．
\end{itemize}

\subsection{题目6}
6．求 \(f(x)=(x-1) x^{\frac{2}{3}}\) 的极小值 \(\_\_\_\_\)．

\subsection{题目7}
7．求极限 \(\lim _{n \rightarrow \infty}\left[\left(n^{2}+1\right)^{\frac{1}{n}}-1\right] \sin \frac{n \pi}{2}=\) \(\_\_\_\_\)．

\subsection{题目8}
8．设函数 \(f(x)\) 是以 \(2 \pi\) 为周期的连续函数，若 \(f(x)\) 在 \([-\pi, \pi]\) 上的傅里叶级数展开式为

\[
\frac{a_{0}}{2}+\sum_{n=1}^{\infty}\left(a_{n} \cos n x+b_{n} \sin n x\right) .
\]

求 \(\frac{1}{\pi} \int_{-\pi}^{\pi} f^{2}(x) \mathrm{d} x=\) \(\_\_\_\_\)．

\subsection{题目9}
9．求 \(\ln \left(1-x+x^{2}\right)\) 在 \(x=0\) 处的幂级数展开式 \(\_\_\_\_\)．

\subsection{题目10}
10．设 \(f(x)\) 在 \(x=0\) 的某邻域内有定义，且 \(\lim _{x \rightarrow 0} \frac{f(x)}{x^{2}}=1\) ，求 \(\lim _{x \rightarrow 0}\left(1+\frac{f(x)}{x}\right)^{\frac{1}{x}}\) ．

\subsection{题目11}
11．求曲面积分

\[
\iint_{\sum} y^{2} z \mathrm{~d} x \mathrm{~d} y+x z \mathrm{~d} y \mathrm{~d} z+x^{2} y \mathrm{~d} z \mathrm{~d} x .
\]

其中曲面 \(\sum\) 为 \(z=x^{2}+y^{2}, x^{2}+y^{2}=1, z=0\) 所围区域的表面，取外侧．

\subsection{题目12}
12．\(f(x, y)=\left\{\begin{array}{ll}\frac{x y\left(x^{2}-y^{2}\right)}{x^{2}+y^{2}}, & (x, y) \neq(0,0), \\ 0, & (x, y)=(0,0) .\end{array}\right.\) 求 \(f_{x y}(0,0)\) 和 \(f_{y x}(0,0)\) ．

\subsection{题目13}
13．证明： \(\operatorname{sgn}(x)=\left\{\begin{array}{ll}-1, & x<0, \\ 0, & x=0, \\ 1, & x>0 .\end{array}\right.\) 在 \([-1,1]\) 上黎曼可积，但在 \([-1,1]\) 上并没有原函数．

\subsection{题目14}
14．已知 \(f(x)\) 在 \([a, b]\) 上连续，证明：\(f(x)\) 在 \([a, b]\) 上有最大值．

\subsection{题目15}
15．判断当 \(\alpha\) 取何值时，函数列 \(f_{n}(x)=\frac{x(\ln n)^{\alpha}}{n^{x}}\) 在 \([0,+\infty)\) 上一致收敛？

\subsection{题目16}
16．解答如下问题：\\
（1）已知级数 \(\sum_{n=1}^{\infty} a_{n}, \sum_{n=1}^{\infty} c_{n}\) 收敛，且 \(a_{n} \leq b_{n} \leq c_{n}\) ，判断 \(\sum_{n=1}^{\infty} b_{n}\) 是否收敛？若收敛，请证明，若不收敛，请给出反例。\\
（2）对于 \(\sum_{n=1}^{\infty} \frac{(-1)^{n-1}}{n^{p+\frac{1}{n}}}\) ，根据 \(p\) 的值讨论该级数何时条件收敛，何时绝对收敛？
\end{document}
